%----------------------------------------------------------------------------------------------
%------------------------------------     2 Determinante    -----------------------------------
%----------------------------------------------------------------------------------------------	
\section{Determinante}
Die Determinante ist eine Kennzahl dafür, ob eine \textbf{quadratische} Matrix singulär oder regulär ist. 
\vspace{0.2cm}

\begin{tabular}{lll}
	Regulär: 	& genau 1 Lösung 			& $\Rightarrow \quad \det(A) \neq 0$ \\
	Singulär: 	& 0 oder $\infty$ Lösungen 	& $\Rightarrow \quad \det(A) = 0$ 	\\
\end{tabular}
		

%----------------------------------------------------------------------------------------------
%------------------------------     2.1 Notation Determinante    ------------------------------
%----------------------------------------------------------------------------------------------	
\subsection{Notation Determinante}

\begin{tabular}{clcl}
	&$A =
	\begin{pmatrix}
		a & b & c \\
		d & e & f \\
		g & h & i 
	\end{pmatrix}$ &&
	$\det(A) = 
	\begin{vmatrix}
		a & b & c \\
		d & e & f \\
		g & h & i 
	\end{vmatrix}$ 
\end{tabular}


%----------------------------------------------------------------------------------------------
%------------------     2.2 Definierende Eigenschaften der Determinante    --------------------
%----------------------------------------------------------------------------------------------
\subsection{Definierende Eigenschaften der Determinante}

\begin{enumerate}
	\item Wenn eine Zeile mit einem Faktor $\lambda$ multipliziert wird, so wird auch $\det(A)$ mit $\lambda$ multipliziert
	\item $\det(A)$ ändert nicht bei blauen Operationen
	\item $\det(I) = 1$
	\item Wenn $A$ singulär ist ($\det(A) = 0$) dann ist auch $A^t$ singulär \\ 
		  $\det(A^t) = 0$
\end{enumerate}


%----------------------------------------------------------------------------------------------
%-----------------     2.3 Geometrische Interpretation der Determinante    --------------------
%----------------------------------------------------------------------------------------------
\subsection{Geometrische Interpretation der Determinante}
		
\subsubsection{Fläche}

Der orientierte Flächeninhalt des Parallelogramms, das von zwei Vektoren $\vec{a}$ und $\vec{b}$ aufgespannt wird, ist die Determinante der beiden Vektoren
\begin{center}
	\includegraphics[width=0.7\columnwidth]{images/flaeche-det.PNG}
\end{center}

		
\subsubsection{Volumen}				    

Das orientierte Volumen eines Parallelepipeds (Spat), das von drei Vektoren $\vec{a}$, $\vec{b}$ und $\vec{c}$ aufgespannt wird, ist die Determinante dieser drei Vektoren.\\
\begin{center}
	\includegraphics[width=0.9\columnwidth]{images/volumen-det.PNG}
\end{center}

\columnbreak


%----------------------------------------------------------------------------------------------
%---------------------     2.4 Allgemeine Berechnung der Determinante    ----------------------
%----------------------------------------------------------------------------------------------
\subsection{Allgemeine Berechnung der Determinante}

Die Determinante kann mithilfe des Gauss-Algorithmus berechnet werden.

Die Determinante ist das \textbf{Produkt} aller Pivot-Elemente.

(Rückwärts-Einsetzen beim Gauss-Verfahren nicht nötig)

$$A =
\begin{pmatrix}
	7 & 49 & 14 \\
	-2 & -6 & 60 \\
	-1 & 0 & 57 
\end{pmatrix}$$

\begin{minipage}[t]{0.1\columnwidth}
\end{minipage}
\begin{minipage}[t]{0.21\columnwidth}
\centering
\renewcommand{\arraystretch}{1.4}
	\begin{tabular}{|c c c|}
		\hline
		$\crd{7}$  & $49$ & $14$ \\
		$-2$ & $-6$ & $60$ \\
		$-1$ & $0$  & $57$ \\
		\hline
	\end{tabular}
\smallskip
Pivot-Elemente:
\end{minipage}
\begin{minipage}[t]{0.05\columnwidth}
\centering
	\begin{tabular}{c}
		\\\\
		$\rightarrow$\\[1.8em]
		$\crd{7}$
	\end{tabular}
\end{minipage}
\begin{minipage}[t]{0.175\columnwidth}
\centering
\renewcommand{\arraystretch}{1.4}
	\begin{tabular}{|c c c|}
		\hline
		$1$ & $7$ & $2$ \\
		$0$ & $\crd{8}$ & $64$ \\
		$0$ & $7$ & $59$ \\
		\hline
	\end{tabular}
\end{minipage}
\begin{minipage}[t]{0.05\columnwidth}
\centering
	\begin{tabular}{c}
		\\\\
		$\rightarrow$\\[1.8em]
		$\crd{8}$
	\end{tabular}
\end{minipage}
\begin{minipage}[t]{0.16\columnwidth}
\centering
\renewcommand{\arraystretch}{1.4}
	\begin{tabular}{|c c c|}
		\hline
		$1$ & $7$ & $2$ \\
		$0$ & $1$ & $8$ \\
		$0$ & $0$ & $\crd{3}$ \\
		\hline
	\end{tabular}
\end{minipage}
\begin{minipage}[t]{0.045\columnwidth}
\centering
	\begin{tabular}{c}
		\\\\
		$\rightarrow$\\[1.8em]
		$\crd{3}$
	\end{tabular}
\end{minipage}
\begin{minipage}[t]{0.2\columnwidth}
\centering
\renewcommand{\arraystretch}{1.4}
	\begin{tabular}{|c c c|}
		\hline
		$1$ & $7$ & $2$ \\
		$0$ & $8$ & $64$ \\
		$0$ & $7$ & $59$ \\
		\hline
	\end{tabular}
\end{minipage}